\documentclass[11pt,a4paper]{article}
\usepackage[T1]{fontenc}
\usepackage{lmodern}
\usepackage[margin=27mm]{geometry}
\usepackage{amsmath,amssymb,amsthm,mathtools}
\usepackage{microtype}
\usepackage{booktabs,array}
\usepackage{enumitem}
\usepackage[hidelinks]{hyperref}
\setlength{\parskip}{5pt plus 1pt}
\setlength{\parindent}{0pt}
\setlength{\emergencystretch}{3em}
\setlist{itemsep=2pt,topsep=3pt}
\raggedbottom
\newtheorem{theorem}{Theorem}[section]
\newtheorem{lemma}[theorem]{Lemma}
\newtheorem{corollary}[theorem]{Corollary}
\newtheorem{proposition}[theorem]{Proposition}
\theoremstyle{definition}
\newtheorem{definition}[theorem]{Definition}
\newtheorem{example}[theorem]{Example}
\theoremstyle{remark}
\newtheorem{remark}[theorem]{Remark}
\newcommand{\Z}{\mathbb{Z}}
\newcommand{\N}{\mathbb{N}}
\newcommand{\Q}{\mathbb{Q}}
\newcommand{\F}{\mathbb{F}}
\newcommand{\ps}[1]{#1[[x]]}
\newcommand{\coeff}[2]{[x^{#1}]\,#2}
\newcommand{\hl}[1]{\textbf{#1}}
\newcommand{\para}[1]{\par\medskip\noindent\textbf{#1}\enspace}
\hypersetup{pdftitle={Residues modulo 8 for a series defined by a self-referential exponential},pdfauthor={Dev D. Goyal},pdfsubject={OEIS A397346}}
\title{Residues modulo $8$ for a series defined by a self-referential exponential\\[6pt]\large A proof of the conjecture in OEIS A397346}
\author{Dev D. Goyal}
\date{October 2026}
\begin{document}
\maketitle
\begin{abstract}
The sequence A397346 of the On-Line Encyclopedia of Integer Sequences (OEIS), contributed by Paul~D.~Hanna, lists the
coefficients of the power series $A(x)=\sum_{n\ge0}a(n)x^n=1+x+4x^2+66x^3+2248x^4+\cdots$ determined by
$A(x)=\exp\bigl(x+\sum_{n\ge2}\frac{2n^2-1}{2n^2}a(n)x^n\bigr)$. The entry conjectures that, from $n=2$ on, the
remainders of $a(n)$ on division by $8$ repeat the block $4,2,0,2$. We prove this. We first show that the equation is
equivalent to a recurrence recorded in the entry, and deduce from it that $a(n)$ is an integer divisible by $n$. For the
quotients $b(n)=a(n)/n$ we prove the stronger statement that $b(n)\equiv6\pmod8$ when $n\equiv3\pmod4$ and
$b(n)\equiv2\pmod8$ otherwise, for $n\ge2$. The results have been checked by computer, and the recurrence and the
congruences are proved formally in the Lean proof assistant.
\end{abstract}
\begin{center}\begin{minipage}{0.88\linewidth}\small
\textbf{How this paper was produced.} This paper reports the result of an AI research effort directed by
Dev~D.~Goyal. ChatGPT (OpenAI), using the GPT-6 model with Pro reasoning and working autonomously from prompts by
Dev~D.~Goyal, found the proof, wrote its Lean~4 formalization, and wrote a first draft of the paper.
Dev~D.~Goyal chose the problems and asked for proofs, formal verification and revisions, but did not supply proof
ideas. Claude Opus~5.5, used with Claude Code, then checked the work independently: it rebuilt the formal proofs from
source and confirmed that they rely only on Lean's three standard axioms, recomputed the sequence from its definition
compared it with the published terms and tested every claim numerically, and read every step of the proofs. It also wrote the present
version of this paper.
Dev~D.~Goyal has verified the results and is responsible for them. The paper has not been peer reviewed.
The section ``About this paper'' at the end gives the details.
\end{minipage}\end{center}

\section{Introduction}

\subsection{The sequence}\label{sec:sequence}
The \emph{On-Line Encyclopedia of Integer Sequences} (OEIS)~\cite{oeisdb} is a public database of integer sequences.
Many entries record \emph{conjectures}: patterns observed in the data that the contributor did not prove. This paper
concerns entry A397346~\cite{oeis}, contributed by Paul~D.~Hanna in 2026, which reads in part:
\begin{quote}\small\ttfamily\raggedright O.g.f. satisfies A(x) = exp( x + Sum\_\{n>=2\} (2*n\^{}2-1) * a(n)*x\^{}n / (2*n\^{}2) ).\par\smallskip Conjecture: a(n) == [4,2,0,2] repeating (mod 8) for n >= 2.\par\smallskip (7) a(n) = 2*n*a(n-1) + n*Sum\_\{k=1..n-2\} (2*(k+1)\^{}2-1)/(k+1) * a(k+1)*a(n-k-1) for n > 1, with a(0) = a(1) = 1.\par\smallskip\hfill\textrm{--- OEIS A397346, Paul D. Hanna, June 2026}\end{quote}

Here ``O.g.f.''\ stands for \emph{ordinary generating function}, the series $A(x)=\sum_{n\ge0}a(n)x^n$. Thus
\begin{equation}\label{eq:source}
A(x)=\exp H(x),\qquad H(x)=x+\sum_{n\ge2}\frac{2n^2-1}{2n^2}\,a(n)\,x^n .
\end{equation}
The equation is self-referential: the unknown coefficients $a(n)$ appear on both sides. The last formula quoted, labelled
(7) in the entry, is a recurrence for $a(n)$; Lemma~\ref{lem:equiv} below proves that it is equivalent to
\eqref{eq:source}. The first terms are 1, 1, 4, 66, 2248, 121690, 9370716, 966649250, \dots{} (from $n=0$).

\subsection{The question}
The conjecture says that from $n=2$ on, $a(n)$ modulo $8$ repeats $4,2,0,2$. The table shows the first terms; the
indices in bold are $n\equiv3\pmod4$. The last row shows the remainders of the quotients $b(n)=a(n)/n$, which follow
a simpler pattern.
\begin{center}\small\setlength{\tabcolsep}{3.2pt}
\begin{tabular}{>{\raggedright\arraybackslash}p{5.4em}>{\centering\arraybackslash}p{1.15em}>{\centering\arraybackslash}p{1.15em}>{\centering\arraybackslash}p{1.15em}>{\centering\arraybackslash}p{1.15em}>{\centering\arraybackslash}p{1.15em}>{\centering\arraybackslash}p{1.15em}>{\centering\arraybackslash}p{1.15em}>{\centering\arraybackslash}p{1.15em}>{\centering\arraybackslash}p{1.15em}>{\centering\arraybackslash}p{1.15em}>{\centering\arraybackslash}p{1.15em}>{\centering\arraybackslash}p{1.15em}>{\centering\arraybackslash}p{1.15em}>{\centering\arraybackslash}p{1.15em}>{\centering\arraybackslash}p{1.15em}>{\centering\arraybackslash}p{1.15em}}
\toprule
$n$ & 2 & \hl{3} & 4 & 5 & 6 & \hl{7} & 8 & 9 & 10 & \hl{11} & 12 & 13 & 14 & \hl{15} & 16 & 17 \\ \midrule
$a(n)\bmod 8$ & 4 & \hl{2} & 0 & 2 & 4 & \hl{2} & 0 & 2 & 4 & \hl{2} & 0 & 2 & 4 & \hl{2} & 0 & 2 \\
$\frac{a(n)}{n}\bmod 8$ & 2 & \hl{6} & 2 & 2 & 2 & \hl{6} & 2 & 2 & 2 & \hl{6} & 2 & 2 & 2 & \hl{6} & 2 & 2 \\
\bottomrule
\end{tabular}
\\[4pt]
\begin{tabular}{>{\raggedright\arraybackslash}p{5.4em}>{\centering\arraybackslash}p{1.15em}>{\centering\arraybackslash}p{1.15em}>{\centering\arraybackslash}p{1.15em}>{\centering\arraybackslash}p{1.15em}>{\centering\arraybackslash}p{1.15em}>{\centering\arraybackslash}p{1.15em}>{\centering\arraybackslash}p{1.15em}>{\centering\arraybackslash}p{1.15em}>{\centering\arraybackslash}p{1.15em}>{\centering\arraybackslash}p{1.15em}>{\centering\arraybackslash}p{1.15em}>{\centering\arraybackslash}p{1.15em}>{\centering\arraybackslash}p{1.15em}>{\centering\arraybackslash}p{1.15em}>{\centering\arraybackslash}p{1.15em}>{\centering\arraybackslash}p{1.15em}}
\toprule
$n$ & 18 & \hl{19} & 20 & 21 & 22 & \hl{23} & 24 & 25 & 26 & \hl{27} & 28 & 29 & 30 & \hl{31} & 32 & 33 \\ \midrule
$a(n)\bmod 8$ & 4 & \hl{2} & 0 & 2 & 4 & \hl{2} & 0 & 2 & 4 & \hl{2} & 0 & 2 & 4 & \hl{2} & 0 & 2 \\
$\frac{a(n)}{n}\bmod 8$ & 2 & \hl{6} & 2 & 2 & 2 & \hl{6} & 2 & 2 & 2 & \hl{6} & 2 & 2 & 2 & \hl{6} & 2 & 2 \\
\bottomrule
\end{tabular}
\end{center}

\subsection{What this paper proves}
\textbf{Main result} (Theorem~\ref{thm:main}). \emph{Equation \eqref{eq:source} has exactly one solution. Its
coefficients $a(n)$ are integers, $n$ divides $a(n)$ for every $n\ge1$, and for $n\ge2$ the quotient $b(n)=a(n)/n$
satisfies}
\[
b(n)\equiv\begin{cases}6\pmod8,&n\equiv3\pmod4,\\2\pmod8,&\text{otherwise}.\end{cases}
\]
Multiplying by $n$ gives the conjecture of the entry (Corollary~\ref{cor:conj}).

\subsection{What was known}
The entry records several equivalent formulas for $A(x)$, including the recurrence quoted above, and an observation that
$a(n)/(2^n\,n!^2)$ appears to tend to a nonzero limit. When the entry was checked on 24~October~2026, the congruence was
still labelled as a conjecture. We are not aware of an earlier proof.

\subsection{Motivation}
Equations in which a generating function is defined through its own coefficients are a recurring theme in the OEIS, and
the coefficients they produce grow very quickly (here roughly like $2^n\,n!^2$). Divisibility and congruence properties
of such numbers are not visible from the definition, which involves the fractions $\frac{2n^2-1}{2n^2}$ and an
exponential. The proof shows how to remove the denominators: first pass to an equivalent recurrence, then divide out the
factor $n$, and only then reduce modulo $8$.

\subsection{How to read this paper}
Section~\ref{sec:background} recalls power series, the formal derivative and the formal exponential.
Section~\ref{sec:results} states the results and Section~\ref{sec:proofs} proves them. Section~\ref{sec:checks}
describes the computer checks and the formal proof, including what is proved in Lean and what only on paper, and
Section~\ref{sec:discussion} discusses the scope of the result.
\section{Background}\label{sec:background}
This section collects the definitions and facts used later. Readers who know formal power series and congruences can skip it and return when needed.
\subsection{Formal power series}
A \emph{formal power series} is an expression
\[
U(x)=u_0+u_1x+u_2x^2+u_3x^3+\cdots ,
\]
determined by its sequence of coefficients $u_0,u_1,u_2,\dots$; see~\cite{niven,stanley}. No number is ever
substituted for $x$, so questions of convergence do not arise. We write $\coeff{n}{U}=u_n$ for the coefficient of $x^n$ in $U$; two series are equal
when all their coefficients are equal.

Series are added coefficient by coefficient and multiplied like polynomials:
\[
\coeff{n}{(UV)}=u_0v_n+u_1v_{n-1}+\cdots+u_nv_0 .
\]
Each coefficient of a product is a finite sum. For example, $(1-x)(1+x+x^2+x^3+\cdots)=1$, so $1+x+x^2+\cdots$ is
the \emph{multiplicative inverse} of $1-x$, written $\frac{1}{1-x}$. More generally, every series $D$ with constant term
$1$ has a unique multiplicative inverse $\frac1D$, whose coefficients are determined one at a time by $D\cdot\frac1D=1$
using only addition, subtraction and multiplication. Whenever a quotient of power series appears in this paper, its
denominator has constant term $1$, so the quotient is again a power series.

A \emph{commutative ring} is a set with an addition, a subtraction and a multiplication that satisfy the usual rules
of arithmetic; examples are $\Z$, $\Q$ and the integers modulo $m$. For a commutative ring $R$ we write $\ps{R}$ for
the set of power series with coefficients in $R$; for instance, writing $\F_2$ for the integers modulo $2$, $\ps{\F_2}$
consists of the series with coefficients modulo $2$. We write
$U(0)$ for the constant term $u_0$ of $U$, and $U=O(x^d)$ when the coefficients of $1,x,\dots,x^{d-1}$ in $U$ are all
zero. Thus $U-V=O(x^d)$ means that $U$ and $V$ have the same coefficient of $x^n$ for every $n<d$. A series with
constant term $0$ and coefficient of $x$ equal to $1$, that is, a series of the form $x+u_2x^2+u_3x^3+\cdots$, is
called \emph{normalized}.

\subsection{The formal derivative and the formal exponential}
The \emph{formal derivative} of $U=\sum_nu_nx^n$ is $U'=\sum_{n\ge1}nu_nx^{n-1}$. It obeys the usual rules
$(U+V)'=U'+V'$ and $(UV)'=U'V+UV'$, which are checked coefficient by coefficient.

If $H$ has rational coefficients and constant term $0$, the \emph{formal exponential}
\[
\exp H=\sum_{r\ge0}\frac{H^r}{r!}
\]
is a well-defined power series: $H^r$ has no terms below degree $r$, so each coefficient is a finite sum. Since
$(H^r)'=rH^{r-1}H'$, differentiating term by term gives $(\exp H)'=H'\exp H$, and $\exp H$ has constant term $1$.

\begin{lemma}\label{lem:ode}
Let $H$ have rational coefficients and constant term $0$. Then $\exp H$ is the only series $F$ with rational coefficients
such that $F(0)=1$ and $F'=FH'$.
\end{lemma}

\begin{proof}
Write $F=\sum_nf_nx^n$ and $H'=\sum_ih_ix^i$. Comparing the coefficients of $x^n$ in $F'=FH'$ gives
$(n+1)f_{n+1}=\sum_{k=0}^{n}f_k\,h_{n-k}$, where the right side involves only $f_0,\dots,f_n$.
So $f_{n+1}$ is determined by the earlier coefficients (we may divide by $n+1$ because the coefficients are rational),
and $f_0=1$ is given. Hence there is at most one such $F$, and $\exp H$ is one.
\end{proof}

\subsection{Remainders and congruences}
For an integer $m\ge2$ we write $a\equiv b\pmod m$ (``$a$ is congruent to $b$ modulo $m$'') when $a-b$ is divisible by
$m$, that is, when $a$ and $b$ leave the same remainder after division by $m$. The remainder of $a$ on division by
$m$ is called its \emph{residue} modulo $m$; it is one of $0,1,\dots,m-1$.

Addition and multiplication are compatible with congruences, so the residues $0,1,\dots,m-1$, added and multiplied
modulo $m$, form a ring, written $\Z/m\Z$. When $m=p$ is a prime we also write $\F_p$. In $\F_2=\{0,1\}$ we have $1+1=0$,
so addition and subtraction coincide.

For power series with integer coefficients, $U\equiv V\pmod m$ means that \emph{every} coefficient of $U-V$ is
divisible by $m$. Equivalently, the two series become equal when every coefficient is replaced by its residue; we call
this \emph{reducing} the series modulo $m$. The key fact is the following.

\emph{Reduction modulo $m$ is compatible with all the operations above.} To reduce a sum, a product, a substitution, or
an inverse, we may reduce the ingredients first and then compute in $\Z/m\Z$. This holds because every
coefficient of the result is obtained from finitely many coefficients of the ingredients by adding, subtracting and
multiplying, and these operations are compatible with remainders.

\section{Results}\label{sec:results}

\begin{theorem}\label{thm:main}
There is exactly one sequence of rational numbers $a(0),a(1),\dots$ satisfying \eqref{eq:source}. Its terms are
integers, $a(0)=a(1)=1$, and $n$ divides $a(n)$ for every $n\ge1$. The quotients $b(n)=a(n)/n$ satisfy, for every $n\ge2$,
\begin{equation}\label{eq:bres}
b(n)\equiv\begin{cases}6\pmod8,&n\equiv3\pmod4,\\2\pmod8,&n\not\equiv3\pmod4.\end{cases}
\end{equation}
\end{theorem}

\begin{corollary}[Hanna's conjecture]\label{cor:conj}
For every $k\ge0$,
\[
a(4k+2)\equiv4,\qquad a(4k+3)\equiv2,\qquad a(4k+4)\equiv0,\qquad a(4k+5)\equiv2\pmod8 .
\]
\end{corollary}

\begin{proof}
By Theorem~\ref{thm:main}, $a(n)=n\,b(n)$. For $n=4k+2$: $(4k+2)\cdot2\equiv4$. For $n=4k+3$: $(4k+3)\cdot6=24k+18\equiv2$.
For $n=4k+4$: $(4k+4)\cdot2\equiv0$. For $n=4k+5$: $(4k+5)\cdot2\equiv2$. All congruences are modulo $8$.
\end{proof}

\section{Proofs}\label{sec:proofs}

\para{The idea in brief.} Differentiating \eqref{eq:source} turns it into a recurrence for $a(n)$ (formula (7) of the
entry), which contains only the fractions $\frac{2j^2-1}{j}$. Writing $a(n)=n\,b(n)$ removes these fractions, and gives a
recurrence for $b(n)$ with integer coefficients. Modulo $8$, every product $b(j)\,b(n-j)$ of two terms with $j,n-j\ge2$ is
$4$, and the recurrence collapses to $b(n)\equiv-b(n-1)+(\text{an explicit correction})$, which can be checked in the four
classes of $n$ modulo $4$.

\begin{lemma}\label{lem:equiv}
A sequence of rational numbers satisfies \eqref{eq:source} if and only if $a(0)=a(1)=1$ and, for every $n\ge2$,
\begin{equation}\label{eq:rec}
a(n)=2n\,a(n-1)+n\sum_{j=2}^{n-1}\frac{2j^2-1}{j}\,a(j)\,a(n-j).
\end{equation}
\end{lemma}

\begin{proof}
Suppose \eqref{eq:source} holds. The constant term gives $a(0)=1$, and differentiating gives $A'=AH'$. Here
$H'=1+\sum_{j\ge2}\frac{2j^2-1}{2j}a(j)x^{j-1}$. The constant term of $A'=AH'$ gives $a(1)=a(0)=1$. For $n\ge2$, the
coefficient of $x^{n-1}$ gives
\begin{equation}\label{eq:coef}
n\,a(n)=a(n-1)+\sum_{j=2}^{n-1}\frac{2j^2-1}{2j}\,a(j)\,a(n-j)+\frac{2n^2-1}{2n}\,a(n),
\end{equation}
where the last term is the summand $j=n$ (with $a(0)=1$). Moving it to the left, the coefficient of $a(n)$ becomes
$n-\frac{2n^2-1}{2n}=\frac1{2n}$; multiplying by $2n$ gives \eqref{eq:rec}.

Conversely, suppose $a(0)=a(1)=1$ and \eqref{eq:rec} holds. Reading the computation backwards gives \eqref{eq:coef} for
every $n\ge2$, and together with $a(1)=a(0)$ this says $A'=AH'$, where $H$ is defined from the sequence as in
\eqref{eq:source}. Since $A(0)=1$, Lemma~\ref{lem:ode} gives $A=\exp H$.
\end{proof}

\begin{proof}[Proof of Theorem~\ref{thm:main}]
\emph{Existence, uniqueness and integrality.} Recurrence \eqref{eq:rec} determines $a(n)$ from $a(0),\dots,a(n-1)$, so it
has exactly one solution, and by Lemma~\ref{lem:equiv} so does \eqref{eq:source}. Define integers $b(n)$ by $b(1)=1$ and,
for $n\ge2$,
\begin{equation}\label{eq:brec}
b(n)=2(n-1)\,b(n-1)+\sum_{j=2}^{n-1}(2j^2-1)(n-j)\,b(j)\,b(n-j).
\end{equation}
Put $a(0)=1$ and $a(n)=n\,b(n)$ for $n\ge1$. Substituting $a(j)=j\,b(j)$ into \eqref{eq:rec} and dividing by $n$ gives
exactly \eqref{eq:brec}, so this sequence satisfies \eqref{eq:rec}. By uniqueness it is the solution, which is therefore
integral with $n\mid a(n)$.

\emph{The residues.} We prove \eqref{eq:bres} by strong induction on $n\ge2$; note $b(2)=2$. Let $n\ge3$. The summand
$j=n-1$ of \eqref{eq:brec} is $(2(n-1)^2-1)\,b(n-1)$, because $b(1)=1$. Combining it with the first term,
\begin{equation}\label{eq:split}
b(n)=\bigl(2n(n-1)-1\bigr)\,b(n-1)+\sum_{j=2}^{n-2}(2j^2-1)(n-j)\,b(j)\,b(n-j).
\end{equation}
In the remaining sum, both $j$ and $n-j$ lie between $2$ and $n-2$, so by induction $b(j)$ and $b(n-j)$ are each $2$ or
$6$ modulo $8$, and their product is $4$ modulo $8$ (as $2\cdot2$, $2\cdot6$ and $6\cdot6$ all are). Since $2j^2-1$ is odd,
each summand is $\equiv4(n-j)\pmod8$. Also $2n(n-1)$ is divisible by $4$ and $b(n-1)$ is even, so the first term is
$\equiv-b(n-1)$. Using $\sum_{j=2}^{n-2}(n-j)=\sum_{i=2}^{n-2}i=\frac{(n-2)(n-1)}2-1$, we obtain
\[
b(n)\equiv-b(n-1)+2(n-2)(n-1)-4\pmod 8 .
\]
(For $n=3$ the sum is empty and the correction $2\cdot1\cdot2-4$ is $0$, so the formula holds there too.) The four cases:
\begin{center}\small\renewcommand{\arraystretch}{1.2}
\begin{tabular}{cccc}
\toprule
$n\bmod4$ & $b(n-1)\bmod8$ (induction) & $2(n-2)(n-1)-4\bmod8$ & $b(n)\bmod8$\\
\midrule
$0$ & $6$ & $0$ & $2$\\
$1$ & $2$ & $4$ & $2$\\
$2$ & $2$ & $4$ & $2$\\
$3$ & $2$ & $0$ & $6$\\
\bottomrule
\end{tabular}
\end{center}
(For $n=3$ the value $b(2)=2$ is used directly.) This is \eqref{eq:bres}, which completes the induction.
\end{proof}

\section{Checking the results}\label{sec:checks}

A mathematical proof can contain mistakes, so the results were checked in two independent ways.

\para{Numerical check.}
We computed $a(n)$ directly from the definition \eqref{eq:source}, with exact rational arithmetic, for $n<61$:
the coefficient of $x^n$ in $\log A$ must equal $\frac{2n^2-1}{2n^2}a(n)$, which determines $a(n)$ from the earlier
terms. This reproduces all 18 terms printed in the OEIS entry, and confirms that the terms are integers. We then computed
$a(n)$ from \eqref{eq:rec} for $n\le400$; the two computations agree for $n<61$, and the computation confirms
$n\mid a(n)$, the residues of Corollary~\ref{cor:conj}, and \eqref{eq:bres} throughout that range.
The computation was done with the PARI/GP computer algebra system, using code that was written separately from the
formal proofs (file \texttt{verification/independent\_checks.gp} in the accompanying repository~\cite{repo}).
A computation can only test finitely many cases, so it is a check on the definitions and on the indexing, not a proof.

\para{Formal proof in Lean.}
Lean~\cite{lean} is a \emph{proof assistant}: a program in which mathematical statements and their proofs are written
in a precise formal language, and every step is checked mechanically down to the basic rules of logic. The theorems of Section~\ref{sec:results} are proved in Lean~4 (version 4.19.0) on top of
Mathlib~\cite{mathlib}, Lean's library of formalized mathematics; the steps proved only on paper are identified below. The statements are collected in the file \texttt{lean/Entries/A397346.lean} of the repository,
in the namespace \texttt{OEIS.A397346}. In Lean the sequence is defined by the recurrence \eqref{eq:rec}, formula (7) of the OEIS entry, and the
identity $A'=AH'$ is proved for it. The last step of Lemma~\ref{lem:equiv}, from $A'=AH'$ to $A=\exp H$
(Lemma~\ref{lem:ode}), is proved in this paper but not in Lean. The proofs are in
\texttt{lean/Families/ExpRecurrence.lean} (namespace \texttt{Families.ExpRecurrence}).
The table lists where each result of this paper is proved.
\begin{center}\small\renewcommand{\arraystretch}{1.25}
\begin{tabular}{>{\raggedright\arraybackslash}p{0.52\linewidth}>{\raggedright\arraybackslash}p{0.4\linewidth}}
\toprule
Statement in this paper & Lean theorem(s)\\
\midrule
The sequence satisfies \eqref{eq:rec} and has the initial values $a(0)=a(1)=1$ & \texttt{canonical\_\allowbreak{}literal\_\allowbreak{}recurrence}, \texttt{displayed\_\allowbreak{}initial\_\allowbreak{}values} \\
\eqref{eq:rec} has exactly one rational solution & \texttt{literal\_\allowbreak{}recurrence\_\allowbreak{}existence\_\allowbreak{}uniqueness} \\
Every rational solution of \eqref{eq:rec} is integral, $n\mid a(n)$, and Corollary~\ref{cor:conj} holds & \texttt{every\_\allowbreak{}source\_\allowbreak{}solution\_\allowbreak{}is\_\allowbreak{}integer\_\allowbreak{}and\_\allowbreak{}has\_\allowbreak{}the\_\allowbreak{}conjectured\_\allowbreak{}residues} \\
Corollary~\ref{cor:conj} & \texttt{original\_\allowbreak{}all\_\allowbreak{}index\_\allowbreak{}conjecture} \\
Theorem~\ref{thm:main}, the residues \eqref{eq:bres} & \texttt{stronger\_\allowbreak{}quotient\_\allowbreak{}congruence} \\
The differential equation $A'=AH'$ of Lemma~\ref{lem:equiv} & \texttt{exact\_\allowbreak{}source\_\allowbreak{}differential\_\allowbreak{}identity} \\
\bottomrule
\end{tabular}
\end{center}
The Lean project was rebuilt from source, and the axioms on which each of these theorems depends were listed with
Lean's \texttt{\#print axioms} command. Every one depends only on the three standard axioms of Lean's logic
(\texttt{propext}, \texttt{Classical.choice}, \texttt{Quot.sound}); no step is left unproved. A formal proof certifies
exactly the statement that Lean checks. Whether that statement says what the OEIS entry means is a separate question,
so the Lean definitions were also compared with the definition in the OEIS entry
(Section~\ref{sec:sequence}).

\section{Discussion}\label{sec:discussion}

\para{Scope.} Theorem~\ref{thm:main} settles the congruence conjecture of the entry. It does not address the observed
limit of $a(n)/(2^n\,n!^2)$, which concerns the size of $a(n)$.

\para{Modulo $16$.} A computer calculation suggests that the pattern continues one level further: for $2\le n<200$, the
remainders of $a(n)$ modulo $16$ repeat $4,2,8,10,12,2,0,2$ with period $8$, and those of $b(n)$ repeat
$2,6,2,2,10,14,2,2$. The argument above would need to keep track of the products $b(j)\,b(n-j)$ modulo $16$, which are no
longer constant.

\para{The role of the factor $n$.} The fractions in \eqref{eq:source} have denominators $2n^2$, and those in
\eqref{eq:rec} have denominators $j$. Both disappear for the quotients $b(n)=a(n)/n$, which satisfy a recurrence with
integer coefficients. The main technical point of the proof is to pass to these quotients before reducing
modulo $8$.

\section*{About this paper}\label{sec:about}
\para{How the work was done.}
This paper is one of 29 papers on conjectures from the OEIS that came out of an AI research effort directed by
Dev~D.~Goyal~\cite{repo}. The roles were as follows.
\begin{itemize}
\item \emph{Dev~D.~Goyal} started the project, chose which problems to pursue, asked for proofs, formal verification and
revisions, and decided what to publish. Dev~D.~Goyal did not supply proof strategies or intermediate arguments.
\item \emph{ChatGPT} (OpenAI), using the GPT-6 model with Pro reasoning, worked autonomously from those prompts. It
searched the OEIS for suitable conjectures, found the proofs, wrote the Lean~4 formalization, and wrote a first draft
of this paper.
\item \emph{Claude Opus~5.5}, used with Claude Code, checked the work independently. It rebuilt the Lean project from source and confirmed that the theorems listed in Section~\ref{sec:checks} rely only on
Lean's standard axioms; compared the Lean definitions with the OEIS definition; recomputed the sequence with separately
written PARI/GP code, compared it with the terms published in the OEIS, and tested every claim of the paper on many
terms; and read every step of the proofs. It then wrote the present version of the paper, expanding the
explanations and, in places, writing out an argument in more detail or by a more elementary route. The theorems of
Section~\ref{sec:results} and the proof strategies are ChatGPT's. The Discussion section, including the short arguments
in it, was written by Claude Opus~5.5; its observations were checked by computer and are not part of the Lean
formalization.
\end{itemize}
Dev~D.~Goyal has verified the results and is responsible for this paper.

\para{Review status.} This paper has not been peer reviewed, and it has not been submitted to a journal or to the OEIS.

\para{Priority.} An OEIS entry that labels a statement as a conjecture shows only that its contributor did not know
a proof; it does not show that no proof existed. During the research only limited searches for earlier work were
made. No claim of priority is made.

\para{Corrections.} Please report any error in the issue tracker of the repository~\cite{repo}, where the Lean sources and the numerical checks are also
available.

\section*{Acknowledgements}
The sequence, its conjecture and the recurrence used here are due to Paul~D.~Hanna~\cite{oeis}.

\begin{thebibliography}{99}\small
\raggedright
\bibitem{oeis} P.~D.~Hanna, entry A397346 (O.g.f.\ satisfies $A(x)=\exp\bigl(x+\sum_{n\ge2}(2n^2-1)a(n)x^n/(2n^2)\bigr)$), \emph{The On-Line Encyclopedia of Integer Sequences}, June 2026. \url{https://oeis.org/A397346}. Text as of October 2026.
\bibitem{oeisdb} N.~J.~A.~Sloane and The OEIS Foundation Inc., \emph{The On-Line Encyclopedia of Integer Sequences}. \url{https://oeis.org}.
\bibitem{stanley} R.~P.~Stanley, \emph{Enumerative Combinatorics}, Volume~2, Cambridge University Press, 1999. (Chapter~5 treats formal power series, composition and compositional inverses.)
\bibitem{niven} I.~Niven, Formal power series, \emph{American Mathematical Monthly} 76 (1969), 871--889.
\bibitem{lean} L.~de~Moura and S.~Ullrich, The Lean~4 theorem prover and programming language, in \emph{Automated Deduction -- CADE 28}, Lecture Notes in Computer Science 12699, Springer, 2021, 625--635.
\bibitem{mathlib} The mathlib Community, The Lean mathematical library, in \emph{Proceedings of the 9th ACM SIGPLAN International Conference on Certified Programs and Proofs (CPP 2020)}, 367--381. Version used: commit \texttt{c44e0c8ee63c}.
\bibitem{repo} D.~D.~Goyal, \emph{Proofs and corrections of OEIS conjectures}, repository, 2026. \url{https://github.com/dvgyl/oeis-a396843-proof}.
\end{thebibliography}
\end{document}
